\section{Ordinary degree and polynomial calculus}\label{sec:algebra}

All polynomial rings in the main argument are over \(\F\).
Degree bounds and accuracy parameters are integers.
Polynomials are ordinary commutative polynomials. In particular, powers are
not silently reduced using Booleanity. The Boolean equations are
\[
 \Bool_x=\{x_j^2-x_j:j\}.
\]
Every displayed generator is an equation with right side zero.

For a finite polynomial system \(\mathcal F\), write
\begin{equation}\label{eq:ns-space}
 \Ideal_B(\mathcal F)
  =\Span_{\F}\{qF:F\in\mathcal F,\ \deg(qF)\le B\}
   \subseteq\Pol_{\le B}.
\end{equation}
The degree constraint is imposed on each generator multiple. Cancellation
between larger-degree multiples does not make them legal in
\(\Ideal_B\).

Ordinary PC derives axioms of \(\mathcal F\), linear combinations of
previous lines, and \(x_jg\) from a previous line \(g\).
A degree-\(B\) derivation has no line of degree exceeding \(B\).
Let \(\Cons_B(\mathcal F)\) be the space of polynomials derivable in this
way. Refutation means deriving one. The space \(\Cons_B\) is closed under
linear combinations, but equality with \(\Ideal_B\) requires a proof.

\begin{lemma}[Multiplying a completed line]\label{lem:pc-product}
If \(g\) has a degree-\(B\) PC derivation and \(q\) is any polynomial,
then \(qg\) has a PC derivation through
\(\max\{B,\deg q+\deg g\}\).
Moreover \(\Ideal_B(\mathcal F)\subseteq\Cons_B(\mathcal F)\).
\end{lemma}
\begin{proof}
Expand \(q\) into monomials. Starting from the already derived line \(g\),
multiply by the variables of each monomial one at a time, and then take
the indicated linear combination. Each new line has degree at most
\(\deg g+\deg q\). The zero case is immediate.
Applying this construction to each legal axiom multiple in
\eqref{eq:ns-space}, then summing, proves the second assertion.
There is no need to multiply the earlier derivation of \(g\).
\end{proof}

\begin{lemma}[Degree-controlled Boolean reduction]\label{lem:boolean}
For every polynomial \(R\) of degree at most \(d\), its squarefree
remainder \(\overline R\) satisfies
\[
 R-\overline R\in\Ideal_d(\Bool_x).
\]
If \(R\) vanishes at every Boolean point, then
\(R\in\Ideal_d(\Bool_x)\). In particular, if \(\deg q\le k\), then
\[
 q^2-q\in\Ideal_{2k}(\Bool_x).
\]
\end{lemma}
\begin{proof}
For \(a\ge2\),
\[
 x^a-x=(x^2-x)(1+x+\cdots+x^{a-2}).
\]
Apply this identity to the powers in each monomial of \(R\), one variable
at a time. Every generator multiple has degree at most the original
monomial degree, so the first assertion follows.

A squarefree polynomial that vanishes on \(\{0,1\}^v\) is zero.
For example, write it as \(A+x_vB\) with \(A,B\) independent of \(x_v\).
Evaluations at \(x_v=0,1\), followed by induction on \(v\), give \(A=B=0\).
This proves the second assertion. Finally, every value of \(q\) at a
Boolean point belongs to \(\F\), where its square equals itself; apply the
second assertion at degree \(2k\).
\end{proof}

We use this reduction lemma only when the required degree is explicitly
available. The high-rank step below instead needs literal membership in
an affine ideal; pointwise vanishing on a finite field does not by itself
supply that stronger conclusion.
